\subsection{实验方法}
设计六组 SysY 程序，分别编写 LLVM IR 与 RISC-V 汇编。
整数和浮点数分别用 i32、float 表示；局部寄存器遵循 SSA，控制流合流使用 phi，
内存变量通过 load/store 访问\cite{llvmref}。

\subsection{编译与验证流程}
使用 Clang/LLVM 检查、优化并编译 IR，链接 getint、putint、putch 完成输入输出。
验证脚本比较参考程序、手写 IR 和手写汇编的标准输出及退出码。
\par\vspace{0.7\baselineskip}\noindent\begin{minipage}[t]{\linewidth}
\begin{lstlisting}[language=bash]
python3 prep/llvm_ir/verify.py
llvm-as sample.ll -o sample.bc
opt -passes=verify -disable-output sample.ll
opt -S '-passes=default<O2>' sample.ll -o sample_O2.ll
llc -mtriple=riscv64-linux-gnu -mattr=+m,+a,+f,+d,+c \
    -target-abi=lp64d sample.ll -o sample_generated_rv.s
\end{lstlisting}
\end{minipage}\par\vspace{0.5\baselineskip}

\subsection{示例设计与特性覆盖}
\begin{table}[htbp]
\centering\small
\begin{tabular}{p{0.23\linewidth}p{0.54\linewidth}r}
\toprule
示例 & 主要覆盖 & 输入数\\
\midrule
arithmetic & 加减乘除与余数、单目正负、常量、块遮蔽、逻辑非转整数 & 6\\
short\_circuit & if/else、逻辑与或短路、全局副作用、自定义函数 & 6\\
loop\_control & while、continue、break、合流与多条回边 & 6\\
functions & 带参函数、递归、整数返回、void 函数及返回 & 5\\
array\_float & 二维数组、部分初始化、数组形参、浮点运算与转换 & 5\\
promote\_demo & 内存形式循环及 mem2reg 对照 & 5\\
\bottomrule
\end{tabular}
\caption{本次手写 LLVM IR 的实验设计}
\end{table}

算术例子输入 $n$，计算 $(n+3)\times2-n/2+n\%3+7$，并分别输出
$-n$、$+n$、$n+!n$。内层同名变量固定为 7，外层 $n$ 在后续输出中仍保持原值，
用于验证遮蔽而非覆盖。负数整除按向零截断计算，余数为
$n-\operatorname{trunc}(n/3)\times3$。
逻辑非先由 icmp 产生 i1，再由 zext 转为 i32 后参与加法。

\subsection{短路控制流与可观测副作用}
短路例子包含 \texttt{n != 0 \&\& 12 / n > 2} 和
\texttt{n == 0 || touch()}。touch 每次调用将全局 hits 加一。
以 n=0 为例，正确程序既不除以零，也不调用 touch，最终 hits=0；
非零输入则调用 touch 一次。

\par\vspace{0.7\baselineskip}\noindent\begin{minipage}[t]{\linewidth}
\begin{lstlisting}
%nonzero = icmp ne i32 %n, 0
br i1 %nonzero, label %and.rhs, label %and.false
and.rhs:
  %q = sdiv i32 12, %n
  %large = icmp sgt i32 %q, 2
  br i1 %large, label %and.true, label %and.false
and.join:
  %and.value = phi i32 [1, %and.true], [0, %and.false]
\end{lstlisting}
\end{minipage}\par\vspace{0.5\baselineskip}
右操作数置于独立基本块中，由分支决定是否执行。
测试 n=-3、0、1、4、5、12，覆盖受保护分母、条件真/假和副作用次数。

\subsection{循环、函数与数据表示}
循环例子逐个处理 1 到 n，跳过 2，在数值超过 5 时退出。
continue 的边直接回到循环头，break 的边直接进入出口。
循环头分别对 i 和 sum 建立 phi，continue 边保持 sum，普通边传入累加后的值。
出口的 phi 对正常条件退出和 break 退出分别提供值。

\par\vspace{0.7\baselineskip}\noindent\begin{minipage}[t]{\linewidth}
\begin{lstlisting}
header:
  %i = phi i32 [0, %entry], [%next, %skip], [%next, %accumulate]
  %sum = phi i32 [0, %entry], [%sum, %skip], [%added, %accumulate]
  %more = icmp slt i32 %i, %n
  br i1 %more, label %body, label %done
done:
  %result = phi i32 [%sum, %header], [%sum, %break.test]
\end{lstlisting}
\end{minipage}\par\vspace{0.5\baselineskip}
LLVM verifier 检查类型、前驱和支配关系，运行测试检查输出。输入 n=2 时只累加 1，n=5 或 n=8 时结果均为 13，
分别触发 continue 与 break。函数例子进一步使用递归阶乘和 void 输出函数，
检查参数、call 和 ret 的对应关系。

数组/浮点例子将 \texttt{int a[2][3]={{1,2,3},{4}}} 表示为二维数组对象，
初始化第二行剩余元素为零。数组实参传地址，按 \texttt{[3 x i32]} 的行类型
通过 getelementptr 寻址，两行和为 10。
浮点过程为 $\operatorname{twice}(1.5n+0.5)=3n+1$，使用 sitofp、fmul、fadd、
fptosi 和 fcmp，既验证 int/float 转换，也验证 float 参数和返回值。


\subsection{mem2reg 观察与结果验证}
另手写一份以 alloca/load/store 保存循环变量的求和程序，使用 LLVM 自带
\texttt{opt -passes=mem2reg} 得到可提升变量的寄存器形式。
转换前为 2 个 alloca、3 个 load、4 个 store、0 个 phi；转换后
这些局部访存操作均消失，形成 2 个 phi。对 n=-3、0、1、5、10
分别运行，转换前后输出一致。
